% ============================================================
%  CV — Mohamed Amine Ziani
%  Compile with pdfLaTeX (works as-is on Overleaf: just upload
%  this file and click "Recompile" — default compiler is fine).
% ============================================================
\documentclass[11pt,a4paper]{article}

% ---------- Page geometry ----------
\usepackage[top=1.7cm,bottom=1.7cm,left=1.9cm,right=1.9cm]{geometry}

% ---------- Fonts & encoding ----------
\usepackage[T1]{fontenc}
\usepackage[utf8]{inputenc}
\usepackage[french]{babel}
\usepackage{tgheros}          % Helvetica-like sans serif, pdfLaTeX-safe
\renewcommand{\familydefault}{\sfdefault}
\usepackage{fontawesome5}     % icons, pdfLaTeX-safe

% ---------- Colors ----------
\usepackage{xcolor}
\definecolor{accent}{HTML}{0E7490}   % teal/cyan, echoes the e-portfolio's cyan accent
\definecolor{darktext}{HTML}{1F2937}
\definecolor{graytext}{HTML}{6B7280}
\color{darktext}

% ---------- Links ----------
\usepackage{hyperref}
\hypersetup{colorlinks=true, linkcolor=accent, urlcolor=accent, citecolor=accent, pdftitle={CV - Mohamed Amine Ziani}}

% ---------- Lists & spacing ----------
\usepackage{enumitem}
\setlist[itemize]{leftmargin=12pt, itemsep=1pt, topsep=2pt, parsep=0pt, label=\textcolor{accent}{\textbullet}}
\setlength{\parindent}{0pt}
\usepackage{setspace}
\setstretch{1.0}
\usepackage{tabularx}

% ---------- Section header macro ----------
\newcommand{\cvsection}[1]{%
  \vspace{7pt}
  {\color{accent}\large\bfseries\scshape #1}\\[1pt]
  {\color{accent}\rule{\linewidth}{0.9pt}}\\[3pt]
}

% ---------- Entry helpers ----------
% \cventry{Title}{Date}{Subtitle/org}{Location}
\newcommand{\cventry}[4]{%
  \noindent\textbf{#1} \hfill {\color{graytext}\small #4}\\
  {\color{accent}\itshape #3} \hfill {\itshape\small #2}\\[2pt]
}

\pagestyle{empty}

\begin{document}

% ============== HEADER ==============
\begin{center}
  {\fontsize{24}{26}\selectfont\bfseries Mohamed Amine Ziani}\\[3pt]
  {\large\color{accent}\bfseries Étudiant Ingénieur ENSEEIHT — Architecture des Ordinateurs et Réseaux (ASR)}\\[6pt]
  {\small
    \faEnvelope\ \href{mailto:medamine.ziani85@gmail.com}{medamine.ziani85@gmail.com}
    \quad\textcolor{graytext}{|}\quad
    \faPhone\ +33\,7\,73\,96\,39\,17
    \quad\textcolor{graytext}{|}\quad
    \faMapMarkerAlt\ Toulouse, France
    \quad\textcolor{graytext}{|}\quad
    \faLinkedin\ \href{https://www.linkedin.com/in/mohamed-amine-ziani}{mohamed-amine-ziani}
    \quad\textcolor{graytext}{|}\quad
    \faGithub\ \href{https://github.com/medamineziani}{medamineziani}
  }
\end{center}
\vspace{2pt}
{\color{accent}\rule{\linewidth}{1.1pt}}

% ============== PROFIL ==============
\cvsection{Profil}
\small
Étudiant ingénieur en 2\ieme{} année à l'ENSEEIHT (filière Architecture des Ordinateurs et Réseaux), issu des CPGE scientifiques au Maroc. Rigoureux et passionné par le développement logiciel, les systèmes/réseaux et l'intelligence artificielle. \textbf{\color{accent}Recherche un stage de 2 mois entre juin et septembre 2027.}
\normalsize

% ============== FORMATION ==============
\cvsection{Formation}
\small
\cventry{ENSEEIHT · INP Toulouse — 2\ieme{} année Ingénieur, Architecture des Ordinateurs et Réseaux (ASR)}{2026 -- Présent}{Architecture des systèmes, réseaux, systèmes d'exploitation, sécurité informatique}{Toulouse, France}

\cventry{ENSEEIHT · INP Toulouse — 1\iere{} année Ingénieur, Sciences du Numérique}{2025 -- 2026}{Algorithmique, programmation, bases de données, réseaux, systèmes d'exploitation}{Toulouse, France}

\cventry{CPGE Scientifique — Mathématiques, Physique, Informatique}{2023 -- 2025}{Baccalauréat Sciences Mathématiques B, Mention Bien (2023)}{Maroc}
\normalsize

% ============== EXPÉRIENCE ==============
\cvsection{Expérience professionnelle}
\small
\cventry{Ingénieur Informatique — AGRIDATA CONSULTING (Stage)}{Juil. 2026 -- Présent, 3 mois}{Deux projets Industrie 4.0 au sein de l'équipe PackOne}{Région d'Agadir, Hybride}
\begin{itemize}
  \item Simulateur IoT industriel en Python/C (sockets TCP/IP, architecture client/serveur) d'une balance et d'un scanner de palettes
  \item Conception de l'infrastructure réseau cible sous Cisco Packet Tracer : segmentation VLAN, routage OSPF, politique QoS pour le trafic de pesée temps réel
  \item PoC Python/OpenCV de contrôle qualité vidéo des palettes (hauteur, nombre de couches de cartons) par caméra en plongée
\end{itemize}

\vspace{4pt}
\cventry{Ingénieur Réseaux IT — FIRST COMPÉTENCE (Stage)}{Juin 2026, 1 mois}{Infrastructure IT et administration réseau}{Agadir, Sur site}
\begin{itemize}
  \item Analyse de l'architecture réseau, adressage IP, VLANs, configuration DHCP/DNS, dépannage et sécurité réseau ; contribution à la documentation de l'infrastructure IT
\end{itemize}
\normalsize

% ============== PROJETS ==============
\cvsection{Projets personnels}
\small
\textbf{Deadlift AI} \hfill \href{https://github.com/medamineziani/deadlift-ai}{\faGithub\ github.com/medamineziani/deadlift-ai}\\
Analyse vidéo temps réel du deadlift : estimation de pose (YOLO11), angles articulaires, trajectoire de la barre, rapport de coaching généré par l'API Claude. \textit{\color{graytext}Python, OpenCV, YOLO, Computer Vision}

\vspace{4pt}
\textbf{PowerForm AI} \hfill \href{https://github.com/medamineziani/powerform-ai}{\faGithub\ github.com/medamineziani/powerform-ai}\\
MVP SaaS complet d'analyse vidéo de powerlifting par IA : authentification/BD Supabase, paiements Stripe, analyse par l'API Claude, déployé sur Vercel. \textit{\color{graytext}Next.js, TypeScript, Supabase, Stripe}

\vspace{4pt}
\textbf{Laboratoire de Cybersécurité}\\
Environnement de pentest personnel (VM Kali Linux + cible Ubuntu, VirtualBox) pour l'auto-formation aux tests d'intrusion : reconnaissance, exploitation, post-exploitation. \textit{\color{graytext}Kali Linux, VirtualBox, Pentest}
\normalsize

% ============== COMPÉTENCES ==============
\cvsection{Compétences techniques}
\small
\begin{tabularx}{\linewidth}{@{}>{\bfseries}l X@{}}
\textbf{Langages} & Python, C, Java, SQL, Ada \\[2pt]
\textbf{Web \& Data} & HTML/CSS, Pandas, NumPy, MATLAB \\[2pt]
\textbf{IA \& Vision} & OpenCV, MediaPipe, YOLO \\[2pt]
\textbf{Réseaux \& Sécurité} & TCP/IP, VLAN, OSPF, QoS, Cisco Packet Tracer, Kali Linux, Pentest \\[2pt]
\textbf{Outils} & Git, Linux, Docker, LaTeX \\[2pt]
\textbf{Méthodes} & POO, Agile, Tests unitaires, Revue de code \\
\end{tabularx}
\normalsize

% ============== RESPONSABILITÉS ==============
\cvsection{Responsabilités}
\small
\cventry{Responsable BDE (Bureau des Élèves) — Classes Préparatoires Scientifiques (CPGE)}{2023 -- 2025}{Animation d'ateliers maths/informatique (10-20 participants), création de supports pédagogiques, coordination d'événements ; également délégué de classe (médiation, suivi des demandes collectives)}{Maroc}
\normalsize

% ============== LANGUES & CENTRES D'INTÉRÊT ==============
\cvsection{Langues \& Centres d'intérêt}
\small
\textbf{Langues :} Français (natif) \textbullet{} Anglais (courant) \textbullet{} Arabe (courant) \textbullet{} Tamazight (courant)\\[2pt]
\textbf{Centres d'intérêt :} Powerlifting \textbullet{} Natation \textbullet{} LeetCode / Codeforces \textbullet{} Veille IA \textbullet{} Open Source
\normalsize

\vfill
\begin{center}
  {\small\color{graytext} Portfolio complet (projets, démos, code) :
  \href{https://github.com/medamineziani/portfolio}{\color{accent}github.com/medamineziani/portfolio}}
\end{center}

\end{document}
